\subtitle{\bf Aula 14 - A Linguagem While}
\begin{frame}[t,plain]
    \titlepage
\end{frame}

\begin{frame}{Objetivos}

    \begin{itemize}
        \item
        Estender a semântica da linguagem com comandos condicionais e de
        repetição.
        \item
        Apresentar a linguagem \emph{TImp}, com declarações de funções e tipos
        registro.
    \end{itemize}
\end{frame}

\begin{frame}{Objetivos}

    \begin{itemize}
        \item
        Definir a semântica big-step para as novas construções.
        \item
        Mostrar como a especificação formal guia a implementação do
        interpretador.
    \end{itemize}
\end{frame}



\begin{frame}{Motivação}

    \begin{itemize}

        \item
        A linguagem \emph{Line} executa apenas sequências de comandos
    \end{itemize}
\end{frame}

\begin{frame}{Motivação}

    \begin{itemize}

        \item
        Programas reais precisam de estruturas de controlo:

        \begin{itemize}

            \item
            \textbf{Condicional}: executar blocos diferentes dependendo de uma
            condição
            \item
            \textbf{Repetição}: executar um bloco enquanto uma condição for
            verdadeira
        \end{itemize}
    \end{itemize}
\end{frame}

\begin{frame}{Sintaxe}

    \[\begin{array}{lcl}
    \mathit{prog} &::=& s^* \\[4pt]
    s &::=& x \mathbin{:=} e \mid \mathbf{read}\; x \mid \mathbf{print}\; e \\
    &\mid& \mathbf{if}\; e\; \mathbf{then}\; \{s^*\}\; \mathbf{else}\; \{s^*\} \\
    &\mid& \mathbf{while}\; e\; \mathbf{do}\; \{s^*\} \\[4pt]
    e &::=& v \mid x \mid e_1 + e_2 \mid e_1 * e_2 \\
    &\mid& e_1 < e_2 \mid e_1 = e_2 \mid e_1 \mathbin{|} e_2 \mid \neg e \\[4pt]
    v &::=& n \mid \mathbf{true} \mid \mathbf{false}
\end{array}\]
\end{frame}

\begin{frame}[fragile]{Tipos em Haskell}

    \begin{lstlisting}[language=Haskell]
        data Value = VInt Int | VBool Bool

        data Exp
        = EValue Value | EVar Var
        | Exp :+: Exp  | Exp :*: Exp
        | Exp :<: Exp  | Exp :=: Exp
        | Exp :|: Exp  | ENot Exp

        data Stmt
        = SAssign Var Exp | SRead Var | SPrint Exp
        | SIf    Exp Block Block
        | SWhile Exp Block

        type Block = [Stmt]
    \end{lstlisting}
\end{frame}


\section{Semântica Small-Step}\label{semuxe2ntica-small-step}

\begin{frame}{Expressões}

    \begin{itemize}

        \item
        Idêntico à linguagem Line.

        \begin{itemize}

            \item
            Exceto por novos operadores, que funcionam de maneira similar.
        \end{itemize}
    \end{itemize}
\end{frame}

\begin{frame}{Condicional}

    \begin{itemize}

        \item
        \textbf{Avaliação da condição:}
        \[\frac{(\sigma, e) \to (\sigma, e')}{\left(\begin{array}{ll}
        \sigma,&\; \mathbf{if}\; e\; \mathbf{then}\; B_1\\
        & \mathbf{else}\; B_2\\
    \end{array}\right) \to \left(\begin{array}{ll}\sigma,& \mathbf{if}\; e'\; \mathbf{then}\; B_1\\
    & \mathbf{else}\; B_2\\
\end{array}\right)}
\]
\end{itemize}
\end{frame}

\begin{frame}{Condicional}

    \begin{itemize}

        \item
        \textbf{Ramo verdadeiro:}
        \[\frac{}{(\sigma,\; \mathbf{if}\; \mathbf{true}\; \mathbf{then}\; B_1\; \mathbf{else}\; B_2) \to (\sigma,\; B_1)}\]
    \end{itemize}
\end{frame}

\begin{frame}{Condicional}

    \begin{itemize}

        \item
        \textbf{Ramo falso:}
    \end{itemize}

    \[\frac{}{(\sigma,\; \mathbf{if}\; \mathbf{false}\; \mathbf{then}\; B_1\; \mathbf{else}\; B_2) \to (\sigma,\; B_2)}
    \]
\end{frame}

\begin{frame}{While}

    \begin{itemize}

        \item
        Caso em que há repetição.
    \end{itemize}

    \[\frac{}{(\sigma,\; \mathbf{while}\; e\; \mathbf{do}\; B) \to
    (\sigma,\; \mathbf{if}\; e\; \mathbf{then}\; (B;\; \mathbf{while}\; e\; \mathbf{do}\; B))}
    \]
\end{frame}

\begin{frame}[fragile]{Exemplo}

    \texttt{while i < 3 do \{ i := i + 1 \}}

    \[(\sigma_0,\; \mathbf{while}\; i < 3\; \mathbf{do}\; \{i \mathbin{:=} i+1\})\]
\end{frame}

\begin{frame}{Exemplo}

    \[\xrightarrow{\text{While}}
    (\sigma_0,\; \mathbf{if}\; i < 3\; \mathbf{then}\; \{i \mathbin{:=} i+1;\; \mathbf{while}\ldots\})\]
\end{frame}

\begin{frame}{Exemplo}

    \[\xrightarrow{\text{IfStep, Var, Lt}}
    (\sigma_0,\; \mathbf{if}\; \mathbf{true}\; \mathbf{then}\; \{i \mathbin{:=} i+1;\; \mathbf{while}\ldots\})\]
\end{frame}

\begin{frame}{Exemplo}

    \[\xrightarrow{\text{IfTrue}}
    (\sigma_0,\; i \mathbin{:=} i+1;\; \mathbf{while}\; i < 3\; \mathbf{do}\; \{i \mathbin{:=} i+1\})\]
\end{frame}

\begin{frame}{Exemplo}

    \[\xrightarrow{\text{Assign}}
    ([i \mapsto 1],\; \mathbf{while}\; i < 3\; \mathbf{do}\; \{i \mathbin{:=} i+1\})\]
\end{frame}


\section{Semântica Big-Step}\label{semuxe2ntica-big-step}

\begin{frame}{Expressões}

    \begin{itemize}

        \item
        Similar à semântica de Line.

        \begin{itemize}

            \item
            Apenas introduz novos operadores
        \end{itemize}
    \end{itemize}
\end{frame}

\begin{frame}{Condicional}

    \[\frac{\sigma \vdash e \Downarrow \mathbf{true} \quad \sigma,\; B_1 \Downarrow \sigma'}
    {\sigma,\; \mathbf{if}\; e\; \mathbf{then}\; B_1\; \mathbf{else}\; B_2 \;\Downarrow\; \sigma'}
    \quad\text{(IfTrue)}\]
\end{frame}

\begin{frame}{Condicional}

    \[\frac{\sigma \vdash e \Downarrow \mathbf{false} \quad \sigma,\; B_2 \Downarrow \sigma'}
    {\sigma,\; \mathbf{if}\; e\; \mathbf{then}\; B_1\; \mathbf{else}\; B_2 \;\Downarrow\; \sigma'}
    \quad\text{(IfFalse)}\]
\end{frame}

\begin{frame}{While}

    \begin{itemize}

        \item
        \textbf{Condição falsa (terminação):}
    \end{itemize}

    \[\frac{\sigma \vdash e \Downarrow \mathbf{false}}
    {\sigma,\; \mathbf{while}\; e\; \mathbf{do}\; B \;\Downarrow\; \sigma}
    \]
\end{frame}

\begin{frame}{While}

    \begin{itemize}

        \item
        \textbf{Condição verdadeira (iteração):}
    \end{itemize}

    \[\frac{\begin{array}{c}\sigma \vdash e \Downarrow \mathbf{true}\\ \sigma,\; B
    \Downarrow \sigma'\\ \sigma',\; \mathbf{while}\; e\; \mathbf{do}\; B \Downarrow
    \sigma'' \end{array}}{\sigma,\; \mathbf{while}\; e\; \mathbf{do}\; B
    \;\Downarrow\; \sigma''}\]
\end{frame}

\begin{frame}{Derivação}

    \[\dfrac{
    \dfrac{
    \dfrac{}{\emptyset \vdash 2 \Downarrow 2}
    \quad
    \dfrac{}{\emptyset \vdash 3 \Downarrow 3}
    }{\emptyset \vdash 2 < 3 \Downarrow \mathbf{true}}
    \quad
    \dfrac{
    \dfrac{}{\emptyset \vdash 10 \Downarrow 10}
    }{\emptyset,\; x \mathbin{:=} 10 \;\Downarrow\; [x \mapsto 10]}
    }{\emptyset,\; \mathbf{if}\; 2<3\; \mathbf{then}\; \{x \mathbin{:=} 10\}\; \mathbf{else}\; \{x \mathbin{:=} 20\} \;\Downarrow\; [x \mapsto 10]}\]
\end{frame}


\section{Implementação em
Haskell}\label{implementauxe7uxe3o-em-haskell}

\begin{frame}[fragile]{Condicional}

    \begin{lstlisting}[language=Haskell]
        interpStmt (SIf e blk1 blk2) = do
        v <- interpExp e
        case v of
        VBool True  -> mapM_ interpStmt blk1
        VBool False -> mapM_ interpStmt blk2
        _ -> throwError $ unlines
        ["Expecting a boolean value, but found:", pretty e]
    \end{lstlisting}
\end{frame}

\begin{frame}[fragile]{Repetição}

    \begin{lstlisting}[language=Haskell]
        interpStmt (SWhile e blk1) = do
        v <- interpExp e
        case v of
        VBool False -> pure ()
        VBool True  -> do
        mapM_ interpStmt blk1
        interpStmt (SWhile e blk1)
        _ -> throwError $ unlines
        ["Expecting a boolean value, but found:", pretty e]
    \end{lstlisting}
\end{frame}


\section{A Linguagem TImp}\label{a-linguagem-timp}

\begin{frame}{Motivação}

    \begin{itemize}

        \item
        \emph{While} cobre atribuição, controlo de fluxo e I/O
        \item
        Mas não permite:

        \begin{itemize}

            \item
            Definir \textbf{abstrações reutilizáveis} (funções)
            \item
            Agrupar \textbf{dados heterogêneos} (registros)
        \end{itemize}
    \end{itemize}
\end{frame}

\begin{frame}{Sintaxe Estendida}

    \[\begin{array}{lcl}
    \mathit{prog}  &::=& \mathit{decl}^*\; s^* \\[4pt]
    \mathit{decl}  &::=& \mathbf{record}\; R\; \{\, \overline{f_i \mathbin{:} \tau_i}\,\} \\
    &\mid& \mathbf{fn}\; f\,(\overline{x_i \mathbin{:} \tau_i}) \mathbin{:} \rho\; \{\, s^*\,\} \\[4pt]
    \rho           &::=& \mathbf{void} \mid \tau \\[4pt]
    \tau           &::=& \mathbf{int} \mid \mathbf{bool} \mid \mathbf{string} \mid R
\end{array}\]
\end{frame}

\begin{frame}{Novos Comandos e Expressões}

    \[\begin{array}{lcl}
    s &::=& \cdots \mid \mathbf{var}\; x \mathbin{:} \tau \mathbin{=} e \\
    &\mid& x.f \mathbin{:=} e \mid \mathbf{return}\; e \mid \mathbf{return} \\
    &\mid& f(e_1,\ldots,e_n) \\[4pt]
    e &::=& \cdots \mid e.f \mid \mathbf{new}\; R\,\{\,\overline{f_i \mathbin{=} e_i}\,\} \\
    &\mid& f(e_1,\ldots,e_n)
\end{array}\]
\end{frame}


\section{Domínios Semânticos}\label{domuxednios-semuxe2nticos}

\begin{frame}{Valores Estendidos}

    O domínio de valores cresce com registros e o valor unitário:

    \[v \;::=\; n \;\mid\; b \;\mid\; s \;\mid\;
    \mathbf{rec}(R,\, \{f_1 \mapsto v_1, \ldots, f_n \mapsto v_n\})
    \;\mid\; \mathbf{unit}\]

    \begin{itemize}

        \item
        Registros têm \textbf{semântica de valor}: toda atribuição e passagem
        de parâmetro copia o registro.
    \end{itemize}
\end{frame}

\begin{frame}{Ambiente de Funções}

    Além de \(\sigma\), a semântica usa um \textbf{ambiente de funções}
    global:

    \[\Phi : \mathit{Name} \rightharpoonup (\overline{x_i},\; B)\]

    \begin{itemize}

        \item
        Construído uma vez a partir das declarações
        \item
        Permanece \textbf{imutável} durante a execução
    \end{itemize}
\end{frame}

\begin{frame}[fragile]{Resultado de Bloco}

    O comando \texttt{return} interrompe a execução normal.
    Modelamos isso com:

    \[r \;::=\; \mathbf{normal} \;\mid\; {\uparrow}\, v\]

    \begin{itemize}

        \item
        \(\mathbf{normal}\): terminação sem retorno
        \item
        \({\uparrow}\, v\): \texttt{return} executado, \(v\)
        propagado ao chamador
    \end{itemize}

    As relações de avaliação tornam-se:

    \[\sigma,\Phi \vdash s \Downarrow (\sigma', r) \qquad \sigma,\Phi \vdash B \Downarrow (\sigma', r)\]
\end{frame}


\section{Semântica de Registros}\label{semuxe2ntica-de-registros}

\begin{frame}{Construção de Registro}

    \[\frac{\sigma,\Phi \vdash e_1 \Downarrow v_1 \quad \cdots}
    {\sigma,\Phi \vdash \mathbf{new}\; R\,\{f_1\!=\!e_1,\ldots\}
    \Downarrow \mathbf{rec}(R,\,\{f_1 \mapsto v_1,\ldots\})}
    \]
\end{frame}

\begin{frame}{Acesso a Campo}

    \[\frac{\sigma,\Phi \vdash e \Downarrow \mathbf{rec}(R,\mathit{flds})
    \quad f \in \mathrm{dom}(\mathit{flds})}
    {\sigma,\Phi \vdash e.f \Downarrow \mathit{flds}(f)}
    \quad\text{(Field)}\]
\end{frame}

\begin{frame}{Atribuição de Campo}

    \[\frac{\sigma(x) = \mathbf{rec}(R,\mathit{flds})
    \quad \sigma,\Phi \vdash e \Downarrow v}
    {\sigma,\Phi \vdash x.f \mathbin{:=} e
    \Downarrow (\sigma[x \mapsto \mathbf{rec}(R,\,\mathit{flds}[f \mapsto v])],\;
    \mathbf{normal})}
    \quad\text{(FieldAssign)}\]

    \begin{itemize}

        \item
        Atualiza apenas o campo \(f\); copia modificada vai para \(\sigma\).
    \end{itemize}
\end{frame}


\section{Semântica de Funções}\label{semuxe2ntica-de-funuxe7uxf5es}

\begin{frame}{Propagação de Retorno}

    \textbf{Bloco vazio:} \[\frac{}
    {\sigma,\Phi \vdash \epsilon \Downarrow (\sigma,\; \mathbf{normal})}
    \quad\text{(BlkNil)}\]

    \textbf{Retorno antecipado:}
    \[\frac{\sigma,\Phi \vdash s \Downarrow (\sigma',\; {\uparrow}\,v)}
    {\sigma,\Phi \vdash s\,;\,B \Downarrow (\sigma',\; {\uparrow}\,v)}
    \quad\text{(BlkRet)}\]
\end{frame}

\begin{frame}{Continuação Normal}

    \[\frac{\sigma,\Phi \vdash s \Downarrow (\sigma',\; \mathbf{normal})
    \quad \sigma',\Phi \vdash B \Downarrow r}
    {\sigma,\Phi \vdash s\,;\,B \Downarrow r}
    \quad\text{(BlkCons)}\]
\end{frame}

\begin{frame}{Comandos de Retorno}

    \textbf{Retorno com valor:} \[\frac{\sigma,\Phi \vdash e \Downarrow v}
    {\sigma,\Phi \vdash \mathbf{return}\; e \Downarrow (\sigma,\; {\uparrow}\,v)}
    \quad\text{(RetVal)}\]

    \textbf{Retorno void:} \[\frac{}
    {\sigma,\Phi \vdash \mathbf{return} \Downarrow (\sigma,\; {\uparrow}\,\mathbf{unit})}
    \quad\text{(RetVoid)}\]
\end{frame}

\begin{frame}{Chamada como Expressão}

    \[\frac{\begin{array}{c}
    \Phi(f) = (\overline{x_i},\; B)\\
    \sigma,\Phi \vdash e_i \Downarrow v_i \;\;(\forall\, i)\\
    [x_1 \mapsto v_1,\ldots,x_n \mapsto v_n],\Phi \vdash B
    \Downarrow (\_,\; {\uparrow}\,v)
\end{array}}
{\sigma,\Phi \vdash f(e_1,\ldots,e_n) \Downarrow v}
\quad\text{(CallExpr)}\]

\begin{itemize}

    \item
    Corpo executado em um \textbf{novo ambiente} com apenas os parâmetros.
\end{itemize}
\end{frame}

\begin{frame}{Chamada como Comando}

    \[\frac{\begin{array}{c}
    \Phi(f) = (\overline{x_i},\; B)\\
    \sigma,\Phi \vdash e_i \Downarrow v_i \;\;(\forall\, i)\\
    [x_1 \mapsto v_1,\ldots,x_n \mapsto v_n],\Phi \vdash B \Downarrow (\_,\; \_)
\end{array}}
{\sigma,\Phi \vdash f(e_1,\ldots,e_n) \Downarrow (\sigma,\; \mathbf{normal})}
\quad\text{(CallStmt)}\]

\begin{itemize}

    \item
    Ambiente de chamada devolvido \textbf{inalterado}.
\end{itemize}
\end{frame}


\section{Implementação em
Haskell}\label{implementauxe7uxe3o-em-haskell-1}

\begin{frame}[fragile]{Estado e Monada}

    \begin{lstlisting}[language=Haskell]
        data Value
        = VInt Int | VBool Bool | VString String
        | VRecord Name (Map Field Value)
        | VUnit

        data InterpState = InterpState
        { varEnv  :: Map Var Value    -- sigma
        , funcEnv :: Map Name FuncDecl  -- Phi
        }

        type InterpM a = StateT InterpState (ExceptT String IO) a
    \end{lstlisting}
\end{frame}

\begin{frame}[fragile]{Execução de Bloco}

    \begin{lstlisting}[language=Haskell]
        execBlock :: [Stmt] -> InterpM (Maybe Value)
        execBlock []     = pure Nothing           -- BlkNil
        execBlock (s:ss) = do
        r <- execStmt s
        case r of
        Just _  -> pure r                     -- BlkRet: propagate
        Nothing -> execBlock ss               -- BlkCons: continue
    \end{lstlisting}
\end{frame}

\begin{frame}[fragile]{Criação de Registro}

    \begin{lstlisting}[language=Haskell]
        evalExp (ENew rname initFields) = do
        vals <- mapM (\(f, e) -> (f,) <$> evalExp e) initFields
        pure (VRecord rname (Map.fromList vals))
    \end{lstlisting}
\end{frame}

\begin{frame}[fragile]{Acesso e Atribuição de Campo}

    \begin{lstlisting}[language=Haskell]
        evalExp (EField e f) = do
        v <- evalExp e
        case v of
        VRecord _ fields ->
        case Map.lookup f fields of
        Just val -> pure val
        Nothing  -> throwError ("Record has no field: " ++ f)
        _ -> throwError "Field access on non-record value"
    \end{lstlisting}
\end{frame}

\begin{frame}[fragile]{Atribuição de Campo}

    \begin{lstlisting}[language=Haskell]
        execStmt (SFieldAssign v f e) = do
        rval <- lookupVar v
        newVal <- evalExp e
        case rval of
        VRecord name fields ->
        let fields' = Map.insert f newVal fields
        in  modify (\st -> st
        { varEnv = Map.insert v (VRecord name fields') (varEnv st) })
        _ -> throwError (v ++ " is not a record")
        pure Nothing
    \end{lstlisting}
\end{frame}

\begin{frame}[fragile]{Retorno e Chamada de Função}

    \begin{lstlisting}[language=Haskell]
        execStmt (SReturn Nothing)  = pure (Just VUnit)   -- RetVoid
        execStmt (SReturn (Just e)) = Just <$> evalExp e  -- RetVal

        callFunc :: Name -> [Value] -> InterpM Value
        callFunc fname argVals = do
        fenv <- gets funcEnv
        case Map.lookup fname fenv of
        Nothing -> throwError ("Undefined function: " ++ fname)
        Just (FuncDecl _ params _ body) -> do
        let bindings = zip [v | Param v _ <- params] argVals
        withFreshEnv bindings $ do
        r <- execBlock body
        pure (fromMaybe VUnit r)
    \end{lstlisting}
\end{frame}

\begin{frame}[fragile]{Ambiente Fresco}

    \begin{lstlisting}[language=Haskell]
        withFreshEnv :: [(Var, Value)] -> InterpM a -> InterpM a
        withFreshEnv bindings m = do
        saved <- gets varEnv
        modify (\st -> st { varEnv = Map.fromList bindings })
        r <- m
        modify (\st -> st { varEnv = saved })
        pure r
    \end{lstlisting}

    \begin{itemize}

        \item
        Salva o ambiente corrente, inicializa apenas com os parâmetros,
        restaura ao terminar.
    \end{itemize}
\end{frame}


\section{Conclusão}\label{conclusuxe3o}

\begin{frame}{Conclusão}

    \begin{itemize}
        \item
        \textbf{Condicional}: regras (IfTrue) e (IfFalse) selecionam o ramo
        pela condição.
        \item
        \textbf{While}: (WhileFalse) termina; (WhileTrue) executa o corpo e
        recomeça.
    \end{itemize}
\end{frame}

\begin{frame}{Conclusão}

    \begin{itemize}

        \item
        \textbf{Registros}: \(\mathbf{rec}(R, \mathit{flds})\) com semântica
        de valor; regras (New), (Field), (FieldAssign).
    \end{itemize}
\end{frame}

\begin{frame}[fragile]{Conclusão}

    \begin{itemize}

        \item
        \textbf{Funções}: resultado de bloco
        \(r \in \{\mathbf{normal}, {\uparrow}\,v\}\) modela
        \texttt{return}; (CallExpr) executa em ambiente
        fresco.
    \end{itemize}
\end{frame}

\begin{frame}{Conclusão}

    \begin{itemize}

        \item
        A especificação formal guia diretamente a implementação em Haskell.
    \end{itemize}
\end{frame}
